\section{Method}

\subsection{Research and Design Approach}

This work follows a design-oriented approach. The target problem is identified from
common academic software project difficulties: unclear requirements, disconnected
tasks and tests, weak evidence, hard-to-review progress, and unclear individual
contribution. DevTrack AI is designed as a system artifact that combines
traceability management, evidence review, GitHub metadata integration, and
AI-assisted reporting.

The method does not claim that AI replaces project management decisions. Instead,
DevTrack AI separates deterministic project logic from AI support. Heuristic and rule-based
logic is used for official status calculation, review gates, artifact relationships,
and evidence scoring. AI is used only for suggestions, structured review summaries,
and report drafts after project data has already been collected.

\subsection{System Artifact Model}

The core model follows the artifact chain introduced in Section~\ref{sec:intro}. A project contains
requirements. Requirements may have use cases and acceptance criteria. Tasks are
linked to requirements and assigned to members. Test cases verify requirements and
may create defects when execution fails. Evidence is uploaded or linked to tasks,
requirements, tests, defects, or code activity. The RTM aggregates these
relationships into a requirement-level health view, and reports are drafted from RTM
snapshots and project activity.

\begin{table}
\centering
\caption{Main DevTrack AI artifacts and their roles}
\begin{tabular}{ll}
\toprule
\textbf{Artifact} & \textbf{Role in the system} \\
\midrule
Requirement & Defines expected system behavior and acceptance criteria. \\
Use Case & Describes actor interaction and requirement scenarios. \\
Task & Represents implementation, testing, design, documentation, or review work. \\
Test Case & Verifies whether a requirement behaves as expected. \\
Bug/Defect & Records failed behavior found during test execution. \\
Evidence & Provides proof of implementation, testing, review, or demo result. \\
RTM & Summarizes traceability and requirement health. \\
Report & Presents weekly progress, risks, blockers, and contribution context. \\
\bottomrule
\end{tabular}
\end{table}

\subsection{Evaluation Direction}

The planned evaluation focuses on traceability completeness and review usefulness.
Student project scenarios can be used to check whether DevTrack AI identifies
requirements without tasks, tasks without evidence, requirements without test cases,
failed tests without defects, unlinked code evidence, and weak completion claims. For
AI-assisted reporting, generated drafts should be compared against structured project
facts to detect unsupported claims. Feedback from leaders, members, and mentors can
also be collected to assess whether the RTM and evidence workflow make project
review clearer.
